\chapter{The Greek Alphabet}\index{Greek alphabet}

Mathematicians and scientists ran out of Latin letters a long time ago, so they borrowed the Greek alphabet. You will see these letters constantly from now on. Here they all are:

\begin{tabular}{c c  l | c c l}
Capital & Lower & Pronounced & Capital & Lower & Pronounced\\
\hline
$A$ & $\alpha$ & Alpha & $N$ & $\nu$ & Nu\\
$B$ & $\beta$ & Beta & $\Xi$ & $\xi$ & Xi ("ksee" or "zy") \\
$\Gamma$ & $\gamma$ & Gamma & $O$ & $o$ & Omicron\\
$\Delta$ & $\delta$ & Delta & $\Pi$ & $\pi$ & Pi\\
$E$ & $\epsilon$ & Epsilon & $P$ & $\rho$ & Rho\\
$Z$ & $\zeta$ & Zeta & $\Sigma$ & $\sigma$ & Sigma\\
$H$ & $\eta$ & Eta & $T$ & $\tau$ & Tau\\
$\Theta$ & $\theta$ & Theta & $\Upsilon$ & $\upsilon$ & Upsilon\\
$I$ & $\iota$ & Iota & $\Phi$ & $\phi$ & Phi\\
$K$ & $\kappa$ & Kappa & $X$ & $\chi$ & Chi ("Kai")\\
$\Lambda$ & $\lambda$ & Lambda & $\Psi$ & $\psi$ & Psi ("Sigh")\\
$M$ & $\mu$ & Mu & $\Omega$ & $\omega$ & Omega
\end{tabular}

\section{Greek Letters in Math and Science}
Some letters show up so often that they have become nicknames for particular ideas:

\begin{tabular}{c l}
Symbol & Common meaning \\
\hline
$\pi$\index{pi} & circle constant, approximately $3.14159$ \\
$\theta$\index{theta} & an angle \\
$\Delta$ & change in a quantity \\
$\Sigma$\index{summation} & sum of many terms \\
$\lambda$\index{wavelength} & wavelength \\
$\mu$ & coefficient of friction, mean, or micro- \\
$\rho$ & density \\
$\omega$\index{angular velocity} & angular velocity \\
$\Omega$\index{ohm} & ohms, a unit of electrical resistance \\
\end{tabular}

A Greek letter can stand for a fixed \emph{constant} or for a \emph{variable} that changes. For example, the area of a circle is

$$A = \pi r^2$$

where $r$ is the radius. Here $\pi$ is a constant (it is always about 3.14159) and $r$ is a variable. Or take this equation from physics:
$$\Delta x = v \Delta t + \frac{1}{2} a \Delta t^2$$

It gives the change in position, $\Delta x$, for something that starts at velocity $v$ and speeds up with constant acceleration $a$ for a time $\Delta t$. The $\Delta$ (``delta'') means ``change in,'' so $\Delta x$ is read ``delta x.'' All four quantities here are variables.

One more, which you will meet again when we study light: the wavelength of a light wave depends on its frequency,
$$\lambda = \frac{c}{f}$$
where $\lambda$ is the wavelength, $c$ is the speed of light, and $f$ is the frequency. Here $\lambda$ (``lambda'') and $f$ are variables, while $c$ is a constant.

\textbf{Worked example: wavelength of a radio wave.} A radio station broadcasts at a frequency of $100$ MHz $= 1.00 \times 10^8$ Hz. The speed of light is $c = 3.00 \times 10^8$ m/s. Using $\lambda = \frac{c}{f}$,
$$\lambda = \frac{3.00 \times 10^8 \text{ m/s}}{1.00 \times 10^8 \text{ Hz}} = 3.00 \text{ m}$$

\textbf{Worked example: distance under constant acceleration.} A cart starts at $v = 2$ m/s and speeds up at $a = 3$ m/s$^2$ for $\Delta t = 4$ s. Using $\Delta x = v\,\Delta t + \frac{1}{2} a\,\Delta t^2$,
$$\Delta x = (2)(4) + \tfrac{1}{2}(3)(4)^2 = 8 + 24 = 32 \text{ m}$$

\begin{Exercise}[title={Using Greek Symbols}, label=greek1]
A wave has a wavelength of $\lambda = 0.5$ m and travels at $c = 3.00 \times 10^8$ m/s. Using $\lambda = \frac{c}{f}$, find its frequency $f$ in hertz.
\end{Exercise}
\begin{Answer}[ref=greek1]
$f = \frac{c}{\lambda} = \frac{3.00 \times 10^8}{0.5} = 6.00 \times 10^8$ Hz.
\end{Answer}

We will always say what each letter means when we use it. Still, it pays to know the alphabet: an equation is much easier to read when you can say its symbols out loud.